\clearpage
\section*{New buildings before 2006 from the Assessor's records}
\textit{October 2, 2026. Agent-drafted work on branch \texttt{explore-alderman-measures}; Jacob asked to extend
construction back to 1999, so that the 2003, 2015 and 2023 ward maps can all be used, and asked that the data be built
and checked before anything is estimated from it.}

Chicago's building permits begin in 2006, so before then the only record of a new building is the Cook County Assessor's.
\path{tasks/build_assessor_new_buildings} finds new buildings in the Assessor's records for 2000--2025 with one set of
rules, set and checked against the permits of 2006--2022 and applied unchanged to earlier years; whether the Assessor
recorded new buildings in the same way before 2006 cannot be checked building by building, only against the records
below. It lists 34,922 buildings: 22,852 houses and 2,677 flats of 2--6 units from the residential record cards, 6,820
condominium buildings, and 2,573 parcels of apartment-class buildings (classes 3xx and 9xx) forming 1,123 developments.

Four sources were added: the Assessor's assessed values of every Chicago parcel and year, 1999--2026
(\path{tasks/download_assessed_values}, 23.6 million rows); the Clerk's parcel maps of 2000--2025
(\path{tasks/download_parcel_history}), which show which old parcels lie under each new one; the City's Energy
Benchmarking, which gives owners' year built for buildings of 50,000 square feet or more; and the Census Bureau's
annual permit counts for Chicago since 1990. The condominium download (\path{tasks/download_condominium_characteristics})
now takes every Chicago building rather than those with a unit recorded as built in 2000 or later; that filter had
dropped most new buildings of 2000--2002 (31, 60 and 66 found, against 255, 349 and 387), because their units' year
built is recorded inconsistently. The construction buildings built from it (\path{tasks/prepare_permit_construction})
are byte for byte unchanged.

A house or flat is new when a record card first assessed in 2000 or later has a recent year built, or a card's year
built jumps to a recent year. A condominium building is new when its units first appear and were built at most five
years before; older units mark a conversion, and the two separate cleanly. An apartment-class building is new when a
parcel becomes improved the year after a vacant or exempt year; when a new parcel appears on land that was cleared, or
carried little building value, in the three years before (judged from the old parcels under it); or when a parcel
enters an apartment class and the Assessor's 2021+ commercial valuations date its building within two years. Those
valuations also set aside existing buildings (dated more than five years before the event) and date buildings built
while their land was exempt, such as the public housing redevelopments that joined the tax rolls together in 2014.
Each building is counted once: a building first assessed on a parcel retired within three years is counted on the
parcels replacing it, with the earlier date, and apartment-class parcels improved in the same year that touch, or
share a valuation, form one development, usually one building assembled from several lots.

Against the permits (\path{tasks/audits/assessor_new_building_checks}), the rules find 96 percent of the permit-linked
houses and flats issued 2006--2018 within five years, 98 percent of condominium buildings and 92 percent of apartment
buildings, nearly all with the right type. Of buildings found in 2012--2022, 93 to 97 percent of houses and flats and
96 percent of condominium buildings have a new-building permit within 100 feet, against 5 to 15 percent of parcels with
no new building; apartment-class buildings reach 80 to 84 percent on their main routes, against 10 percent. Before
2006, the rules find 81 percent of the surviving apartment buildings the valuations date to 1999--2005 (90 percent for
2006--2020), and a condominium or apartment building lies at 60 percent of the large Benchmarking buildings built then
(54 percent later). Against the Census counts two years before, yearly totals move together (correlations 0.83 for
houses, 0.94 for buildings of five units or more, 0.97 for small multifamily), and large buildings run 0.92 times the
Census before 2006 and 0.94 times after; houses run 1.6 times before and 1.4 times after.

Two hand checks of 103 and 43 sampled buildings (\path{tasks/audits/assessor_new_building_checks/adjudication}) shaped
the rules: the first found buildings counted twice, dates off by one to seventeen years, existing buildings passing as
new and two routes without signal; the second found three narrower problems, each fixed by a rule applied to all
cases. A rule to drop house rebuilds with no rise in value was tested and not adopted, since 79 percent of them have a
permit on the parcel.

\textit{Open.} The routes that rely on the 2021+ valuations see only buildings still standing in 2021, which probably
explains part of the lower recall of early apartment buildings. Separate small buildings on touching parcels built in
the same year count as one development. About ten apartment-class parcels a year carry near-zero values; some are
subsidized buildings and some are not, and no value floor was set. Houses run somewhat higher relative to the Census
before 2006 than after.
